\documentclass{article}
\usepackage[T1]{fontenc}
\usepackage{lmodern}
\usepackage{graphicx}
\usepackage{amsmath,amssymb}
\usepackage[a4paper,margin=2cm]{geometry}
\usepackage[absolute,overlay]{textpos}
\setlength{\TPHorizModule}{1mm}
\setlength{\TPVertModule}{1mm}
\usepackage[indonesian]{babel}
\usepackage{hyperref}
\setlength{\emergencystretch}{2em}
% unit: Halaman 21

\begin{document}

% === Halaman 21 (python/native) ===

void dekripsi()

\{

  string plainteks, cipherteks;

  int i, k;

  char c;

  cout << "Ketikkan cipherteks: ";

  cin.ignore();getline (cin, cipherteks);

  cout << "Masukkan jumlah pergesaran (0-25): ";

  cin >> k;

  plainteks = "";  // inisialisasi plainteks dengan null string 

  for (i=0; i < cipherteks.length(); i++) \{

      c = cipherteks[i];

      if (isalpha(c)) \{    //hanya memproses alfabet

         c = toupper(c);   // ubah karakter ke huruf besar 

         c = c - 65;       // kodekan huruf ke angka 0 s/d 25 

         if (c - k < 0)    // kasus pembagian bilangan negatif

            c = 26 + (c - k);

         else

            c = (c - k) \% 26;       

         c = c + 65;       // kodekan kembali ke huruf semula  

         c = tolower(c); // plainteks dinyatakan sebagai huruf kecil        

      \}      

      plainteks = plainteks + c;  // sambungkan ke plainteks 

  \}

  cout << "Plainteks: " << plainteks << endl;  // cetak plainteks

\}

21

\end{document}
